\documentclass[11pt]{article}
\usepackage{amsmath,amssymb,hyperref}
\title{Two valid custody laws need not share one reusable summary}
\date{25 November 2026}
\begin{document}
\maketitle
\paragraph{What was retained.}
The established \href{https://github.com/eternal-november/eos-research-programmes/pull/391}{fixed-law deadline profile} gives an exact finite continuation test for a stream with fixed carried receipt expiries and declared future mechanical gaps. Use the signer and tag protocol of \href{https://github.com/eternal-november/eos-research-programmes/pull/390}{the shared port model}: a trusted station clock, one physical transfer and mint per slot, alternating signers, two tags per signer in this instance, and an old copy that can persist until one tick before its reuse gap. Before each transfer the current mechanical gap and the same-tag expiry must both be satisfied. Retain the next signer and tag, the last mint time, and a separate exact profile for every law in a fixed menu. Ask if those data suffice when the same authority may lawfully amend one of those laws \,\emph{between} transfers. Neither the amendment nor a new timeout retires an old copy. No independent mission is modeled.
\paragraph{The shortest collision.}
Consider the two successive actual old mint sequences \((-9,-8,-5,-2)\) and \((-9,-7,-5,-2)\), with unit mechanical gaps in the old operation, and reuse gap 10. A past starts with an empty port and proceeds in the protocol's normal alternating tags. The two histories differ only in the second old mint, whose expiry is 2 or 3. With next-signer mechanical gaps \(L=(1,2)\) or \(M=(4,1)\), both histories have respectively the same exact fixed-law earliest prefix:
\[
 L:\;(1,3,5,8),\qquad M:\;(2,3,7,8).
\]
Therefore they have the same full continuation languages under each law held fixed; this is stronger than indistinguishability under only the old law. Now use \(L\) through slot 1 and \(M\) from slot 2, with common verified authority to do so. The first transfer can take place at 1 under \(L\). Under \(M\) the second needs just one mechanical slot, but cannot reuse its signer's old tag before 2 in the first past or 3 in the second. The word \((1,2)\) is lawful only in the first past. Each one-transfer word has the same legality, so this is a shortest distinguishing continuation. Two transfers by slot 2 force this word, as the first possible transfer is at 1 and the port is serial. The result is about a mutable law, not a stronger pure-retirement example of \href{https://github.com/eternal-november/eos-research-programmes/pull/384}{the earlier signature-profile study}.
\paragraph{The exact substitute for an announced calendar.}
Fix a finite horizon \([0,H)\) and a publicly known, authorized function \(d_i(t)>0\) specifying the mechanical gap for transfer index \(i\) at slot \(t\). No new event changes the calendar. Let \(b\) be the last mint, \(K\) the tag cycle, \(e_i\) the fixed old same-tag expiry for \(i<K\), and \(W\) the future reuse gap. Carry the last mint separately. Starting at \(p_{-1}=b\), define
\[
 p_i=\min\{t\in[0,H): t\ge e_i,\ t-p_{i-1}\ge d_i(t)\}\quad(0\leq i<K),
\]
with empty minimum \(H\), which propagates. A previous earlier slot never hurts feasibility of a specified next slot, even when the gap changes with that slot. Any legal prefix therefore has \(t_i\ge p_i\) by induction. Every attainable \(p_i\) is realized by its own greedy prefix. Conversely, supplement the actual chosen gap tests with \(t_i\ge p_i\) in place of the old expiry test, and retain the actual recurrence \(t_i-t_{i-K}\ge W\) for \(i\ge K\). Since \(p_i\ge e_i\) whenever attainable, all tests are then sound and complete. Comparing histories with the same last mint, phase, calendar and recurrence, the clipped prefix is exact for the whole language: if one coordinate differs, the earlier witness prefix satisfies one history and cannot satisfy the other. For a calendar of \(S\) sorted constant intervals, the next candidate within an interval is the maximum of interval start, expiry and previous mint plus current gap. An interval exhausted at this step is never revisited. This gives an \(O(K+S)\) sweep on the normalized interval table, not a speed claim against an equally factored direct solver.
\par
This does not justify plugging in a calendar not yet chosen or using a predicted amendment as authority. A proof required for a family of calendars must cover each admitted history or retain what its active checks consume. The two fixed entries alone are insufficient. For the stated two-past, two-transfer decision it suffices to certify the one bit \(e_1\leq2\) before the second transfer decision. Without evidence, the exact direct history-fibre dispatcher and the checked rule both withhold the deadline guarantee. With timely qualified evidence, both recover the feasible arm. A raw mint timestamp known only to an evaluator supplies no new runtime evidence.
\paragraph{Checks.}
The original selected port-successor implementation and a separate absolute-time checker agree on every subset of the 12-slot horizon under both fixed calendars and the switch. The fixed languages have 321 and 126 words respectively in both pasts; the switched ones have 192 and 171. The new profile implementation also agrees on 61,320 additional comparison words, including orders with rising and falling gaps. The three Lean lemmas check the monotone predecessor, replacement step, and two-transfer contradiction; protocol and authority interpretation are not formalized. The direct same-information comparison ties. This is a sharp proof-reuse boundary, not a stronger deployment or a contract acceptance claim.
\end{document}
